% chapters/17/evaluasi.tex
\bagian{Asesmen dan Instrumen Penilaian}
\label{sec:asesmen-akhir}

\begin{asesmenbox}
\textbf{Kedudukan asesmen bab ini.} Bab ini tidak menuntaskan Sub-CPMK baru.
Ia menjadi penutup rangkaian, dan asesmennya berupa \textbf{Ujian Akhir
Semester} pada minggu keenam belas, beserta penilaian \textbf{laporan
pengembangan} dan \textbf{lembar refleksi} yang menemaninya.

\vspace{0.4em}
UAS mengukur \CPMK{3} dan \CPMK{4}, sesuai penilaian pada RPS. Rubrik lengkap
Proyek Akhir disajikan pada Lampiran~A.
\end{asesmenbox}

\subsection*{Instrumen 1: Kisi-kisi Ujian Akhir Semester}

\begin{table}[htbp]
\centering
\caption{Kisi-kisi Ujian Akhir Semester beserta CPMK yang diukur}
\label{tab:kisi-uas}
{\setlength{\tabcolsep}{4pt}\footnotesize%
\begin{tabularx}{\textwidth}{@{}C{1.7cm} L{4.4cm} C{1.8cm} Y@{}}
\toprule
\headrow \textbf{CPMK} & \textbf{Pokok yang diujikan} & \textbf{Bobot} & \textbf{Bentuk soal} \\
\midrule
\CPMK{1} & \istilah{Game loop} dan \istilah{delta time}; gerak dan vektor;
fisika sederhana; tabrakan &
20\% & Menjelaskan, menghitung, dan memperbaiki potongan kode Pygame. \\
\addlinespace
\CPMK{2} & Perilaku pengejar; mesin keadaan; A* dan heuristiknya &
20\% & Menelusuri keadaan, memilih heuristik, memperbaiki kasus tepi. \\
\addlinespace
\CPMK{3} & \istilah{GameObject} dan komponen; \istilah{Rigidbody} dan
\istilah{Collider}; siklus hidup \code{MonoBehaviour} &
30\% & Menganalisis gejala fisika, dan menentukan letak kesalahannya. \\
\addlinespace
\CPMK{4} & Impor aset; \istilah{Animator}; antarmuka; bunyi; pengujian
prototipe &
30\% & Menganalisis masalah integrasi, dan menyusun rencana perbaikan. \\
\bottomrule
\end{tabularx}}
\end{table}

Separuh bobot UAS berada pada \CPMK{3} dan \CPMK{4}, sebab
UAS adalah satu-satunya instrumen tulis yang mengukur keduanya. Praktikum
mengukur keduanya juga, tetapi melalui artefak, sehingga kemampuan
menjelaskan alasannya tidak terukur di sana.

\subsection*{Instrumen 2: Rubrik Laporan Pengembangan}

\begin{table}[htbp]
\centering
\caption{Rubrik penilaian laporan pengembangan dan refleksi}
\label{tab:rubrik-laporan}
{\setlength{\tabcolsep}{4pt}\footnotesize%
\begin{tabularx}{\textwidth}{@{}L{3.4cm} C{1.3cm} L{3.0cm} Y@{}}
\toprule
\headrow \textbf{Kriteria} & \textbf{Bobot} & \textbf{Bukti yang diperiksa} & \textbf{Deskripsi capaian} \\
\midrule
Kelengkapan dokumen & 20\% & Daftar isi dan lampiran &
Memuat ringkasan desain, daftar aset, lembar \istilah{playtest}, dan
\istilah{log bug} beserta status penanganannya. \\
\addlinespace
Kejelasan laporan perbaikan & 30\% & Bagian perbaikan &
Setiap perbaikan disertai keadaan sebelum dan sesudah, beserta buktinya. \\
\addlinespace
Kesesuaian GDD dan \istilah{build} & 20\% & Tabel kesesuaian &
Perbedaan antara GDD dan \istilah{build} dicantumkan dengan jujur, beserta
alasannya. \\
\addlinespace
Refleksi pemilihan teknologi & 30\% & Lembar refleksi &
Setiap pernyataan memuat keadaan konkret, tindakan, dan hasil; ketiga tinjauan
pada RPS dijawab secara terpisah. \\
\bottomrule
\end{tabularx}}
\end{table}

\subsection*{Instrumen 3: Soal Uraian}

Jawablah dengan singkat dan berdasarkan bukti. Kunci jawaban beserta CPMK yang
diukur disajikan pada Tabel~\ref{tab:kunci-uraian-ch17}.

\begin{enumerate}
  \item Sebuah CPMK memiliki bobot 25 persen, dan CPL yang menaunginya
        memperoleh 57 persen. Jelaskan bagaimana kedua angka itu dihitung.
  \item Mengapa penguasaan setiap bab belum tentu menghasilkan permainan yang
        dapat dimainkan? Berikan satu contoh dari pengalaman Anda.
  \item Pemain tersangkut di sudut dinding pada level Unity. Jelaskan lapisan
        yang perlu diperiksa, berurutan dari bawah.
  \item Sebuah kelompok melaporkan bahwa seluruh kriteria rubrik terpenuhi,
        tetapi penguji menemukan \istilah{soft-lock}. Kriteria mana yang
        seharusnya menangkapnya, dan mengapa lolos?
  \item Tuliskan satu rencana belajar lanjutan beserta langkah pertama yang
        dapat dimulai dalam sebulan.
\end{enumerate}

\begin{table}[htbp]
\centering
\caption{Kunci soal uraian Bab XVII beserta CPMK yang diukur}
\label{tab:kunci-uraian-ch17}
{\setlength{\tabcolsep}{4pt}\footnotesize%
\begin{tabularx}{\textwidth}{@{}C{0.8cm} C{1.7cm} Y@{}}
\toprule
\headrow \textbf{No} & \textbf{CPMK} & \textbf{Pokok jawaban} \\
\midrule
1 & \CPMK{3}, \CPMK{4} &
Bobot CPMK diperoleh dengan menjumlahkan sumbangan setiap komponen kepadanya
pada Tabel~\ref{tab:matriks-bobot}; bobot CPL diperoleh dengan menjumlahkan
bobot CPMK penyusunnya, yaitu $25 + 32 = 57$ persen untuk \CPL{03}. \\
\addlinespace
2 & \CPMK{3}, \CPMK{4} &
Sebab setiap bab melatih satu kemampuan pada keadaan yang sudah disiapkan,
sedangkan permainan yang utuh menuntut kemampuan menggabungkan. Contoh yang
dapat diterima: skor tidak bertambah karena tanda bendanya salah, walaupun
antarmuka dan skripnya benar. \\
\addlinespace
3 & \CPMK{3} &
Tabrakan dan bentuk batas benda terlebih dahulu, lalu cara geraknya dijalankan
pada \code{FixedUpdate}, lalu penggunaan \istilah{delta time}. Penelusuran
bergerak ke bawah sebab lapisan atas bergantung pada lapisan bawah. \\
\addlinespace
4 & \CPMK{4} &
Kriteria ``fungsionalitas mekanik inti dan level'', yang mensyaratkan level
dapat dimainkan dari awal sampai selesai tanpa penghambat. Ia lolos karena
rubrik dinilai dari satu kali penyelesaian, sedangkan \istilah{soft-lock} hanya
dialami pemain yang berada di tempat yang salah. \\
\addlinespace
5 & \CPMK{4} &
Jawaban dinilai dari kekonkretan: satu atau dua kompetensi yang nyata, beserta
satu langkah pertama yang dapat diperiksa waktunya. \\
\bottomrule
\end{tabularx}}
\end{table}

\begin{catatabox}
Soal nomor empat sengaja dipakai sebagai penutup. Ia menuntut Anda menilai
rubriknya sendiri, dan itulah tingkat tertinggi pada ranah C5 yang dinilai
pada mata kuliah ini: bukan sekadar memakai alat penilaian, melainkan
mengetahui batas kemampuannya.
\end{catatabox}
